\documentclass[11pt]{article}

\usepackage[utf8]{inputenc}
\usepackage[T1]{fontenc}
\usepackage{lmodern}
\usepackage{enumitem}
\usepackage[left=2.5cm,right=2.5cm]{geometry}
\usepackage{xcolor}
\definecolor{navy}{RGB}{17,85,170}
\usepackage{amsmath}
\usepackage{graphicx}
\usepackage{booktabs}
\usepackage{subcaption}
\usepackage[capposition=top]{floatrow}
\usepackage{natbib}
\usepackage{currfile}
\usepackage{xr}
\usepackage{xr-hyper}
\usepackage{hyperref}
\externaldocument[][nocite]{main}
\externaldocument[][nocite]{referee3_response}
\newcommand{\comment}[1]{\begingroup #1\endgroup}
\newcommand{\answer}[1]{\begingroup\color{navy}#1\endgroup}
\providecommand{\sym}[1]{#1}

\author{}
\title{Response to the Editor (EJ MS20241010-2)}
\date{}

\begin{document}
%\maketitle

\begin{center}
   \LARGE{Response to Editor (EJ MS20241010-2)}
\end{center}

\vspace{0.5em}
% \noindent \answer{Answer TBW.}\par

% \section*{Roadmap on main empirical changes}

\answer{Answer TBW.

%\begin{enumerate}[leftmargin=*]
%\item \textbf{...}
%\end{enumerate}
}

\par
\vspace{0.5em}


\section*{Editor's letter}


\comment{I went back to the original three referees. I have heard back from R2 and R3; R1 has been unresponsive (note that this was the most positive referee in the first round). I have decided to move ahead without their report.}\par

\comment{Both R2 and R3 are ready to sign off. I have also had a chance to read your paper and your responses (including to R1's comments), as well as the referees' reports. I would like to thank you for submitting a responsive revision and for making a serious attempt to address the concerns that were raised in the previous round. I agree that the paper is substantially improved.}\par

\answer{Thanks for your feedback. We are extremely happy that you assessed the revision and acknowledge that the paper is substantially improved. We also appreciate the time you and the referees took in reviewing our work.}\\


\comment{R3 raises several minor points. I think addressing them (or at least some of them) will strengthen the paper. I would therefore like you to submit one final revised version of the paper. I do \textbf{NOT} plan to go back to the referees, and only I will review the revision. However, I would like you to submit a point-by-point response to R3's comments (again, only I will review them). As always, if you find a comment useful but ultimately decide not to implement it, please explain your reasoning in your response. This is your paper and I am hesitant to micromanage the revision -- it is fine if you decide not to pursue a certain point, but I would like to understand why.}\par

\answer{Even though you're not planning on sending the letter back to the referees, we wrote a full letter addressed to referee 3. We responded to each point, either explaining what we did to implement the point raised or explaining why we decided not to do it.}\\


\comment{I hope to conditionally accept the paper in the next round. Just so that you know what is coming next: after the conditional acceptance, the next step will be to provide the codes and data for replication to the editorial office. I realize there are restrictions in your case -- the data editor will work around those in consultation with you.}\par

\answer{We spend a considerable amount of time making sure the replication package worked correctly. In the process, some numbers moved very marginally, mostly because the underlying data sources have been somewhat updated by the statistical office. At the end of this letter, we comment on the main discrepancies, but none of them are very substantial and they simply reflect the fact that the raw inputs into our work have been updated by the French statistical institute (Insee). }\\


\comment{Please also note that, once the paper is accepted, it will be checked using our anti-plagiarism software. It is therefore advisable to check the accuracy of all references at this stage.}\par

\answer{We made sure all references were properly cited.}\par


%\comment{Thank you for giving us the opportunity to consider your work. I look forward to seeing your revised paper soon (and, subsequently, the accepted paper in print). Please feel free to contact me if you have any questions.}\par
%\answer{Answer TBW.}\par



%\bibliographystyle{aer}
%\bibliography{_REF}
%\bibliography{_REF}

% here claude : main changes of the paper in response to R3 comments. 

\section*{Main changes in estimates between the two versions}

\answer{All results in the revised version come from a complete run of our replication package on the updated source files. Most estimates change only in the second or third decimal, and the qualitative conclusions of the paper are unchanged. We list the main changes below; numbers from the previous version are in parentheses.

\begin{itemize}[leftmargin=*]
\item \textbf{Samples.} The LE sample now includes 21,132 displaced workers (20,897) and the mass-layoff sample 17,946 (17,949). The number of observations available for firm-level outcomes changes more: for instance, it increases by 29\% for productivity in the LE sample (Appendix Table~\ref{apptab:evstudy_firm_outcomes}).

\item \textbf{Employment, earnings and hours.} These estimates are virtually unchanged. One year after displacement, the employment probability of displaced workers is 31.5 percentage points lower (31.7) and their annual earnings are 13,687 euros lower (13,883). Five years after displacement, the losses are $-0.318$ log points for earnings ($-0.304$) and $-0.231$ log points for hours ($-0.223$) (Table~\ref{tab:evstudytable_body}, panel a).

\item \textbf{Hourly wages and firm wage premia.} The hourly-wage loss at $t^D+5$ is $-0.101$ ($-0.098$). The loss in firm wage premia is smaller than in the previous version, mostly in the first years after displacement: $-0.022$ ($-0.047$) at $t^D+1$ and $-0.037$ ($-0.042$) at $t^D+5$. Firm wage premia therefore account for 37\% of the hourly-wage loss at $t^D+5$ (43\%) (Table~\ref{tab:evstudytable_body}, column 7, and Figure~\ref{fig:fe_current_event}). This share is 63\% in the mass-layoff sample (66\%) and 59\% with the unrestricted AKM premium (61\%).

\item \textbf{Productivity and labor share of destination firms.} At $t^D+5$, destination firms are 0.134 log points more productive (0.139) and their labor share is 0.074 log points lower (0.068) (Table~\ref{tab:evstudytable_body}, panel b). The within-sector component accounts for 55\% of the productivity gain (70\%) (Figure~\ref{fig:prod_res_within_btw}).

\item \textbf{Collective agreements.} In the previous version, outcomes related to collective agreements were only measured for firms with more than 100 employees. We now lower this threshold to 50 employees. This makes the sample more representative and aligns it with French labor law: firms with 50 employees or more must set up a works council, and unions can appoint union delegates, who negotiate collective agreements with the employer. The number of observations rises from 74,011 to 89,202 (Appendix Table~\ref{apptab:evstudy_firm_outcomes}). With this broader sample, displaced workers are 5.9 percentage points less likely to work for a firm that signed a wage agreement at $t^D+4$ (7.8) (Appendix Table~\ref{tab:nego:LEnego_acc_result}).

\item \textbf{Professional elections.} Turnout in destination firms is 8.9 percentage points lower at $t^D+5$. The previous version reported two different values for this coefficient, $-0.070$ in Table~\ref{tab:evstudytable_body} and $-0.108$ in Appendix Table~\ref{tab:nego:LEnego_ele_result}; both tables now report the same estimate.
\end{itemize}}\par 

\end{document}
